\documentclass{amsart}
\usepackage{amsmath,amssymb,amsthm,hyperref}
\usepackage{bm}

\theoremstyle{definition}
\newtheorem{definition}{Definition}
\newtheorem{assumption}{Assumption}
\newtheorem{model}{Model}
\newtheorem{remark}{Remark}
\theoremstyle{plain}
\newtheorem{proposition}{Proposition}
\newtheorem{theorem}{Theorem}
\newtheorem{corollary}{Corollary}
\newtheorem{lemma}{Lemma}

\begin{document}

\title{An Integrated Structural--Interpersonal Framework for Team Learning Behavior}
\maketitle

\begin{abstract}
We formulate an integrated framework that jointly models team \emph{structural features}
and \emph{shared interpersonal beliefs} (psychological safety) as determinants of team
\emph{learning behavior} and downstream \emph{performance}. The framework is a recursive
structural causal model (SCM) on a directed acyclic graph (DAG). We are careful to
distinguish three logically separate layers and to prove a precise statement at each:
(a) a \emph{representation/justification} layer, where we list qualitative desiderata read
off the substantive prompt and prove that the proposed model is the (essentially unique)
minimal linear recursive model satisfying them (Theorem~\ref{thm:rep}); (b) a
\emph{statistical (regression) identification} layer, under a disturbance-orthogonality
assumption, giving population least-squares recovery of all coefficients
(Theorem~\ref{thm:ident}); and (c) a \emph{causal} layer, under an additional
sequential-ignorability assumption, under which conditional means equal interventional
means, so the coefficients carry causal meaning (Theorem~\ref{thm:causal}). We then prove an
exact \emph{four-way effect decomposition} matching the canonical
controlled-direct/reference-interaction/mediated-interaction/pure-indirect terms of
VanderWeele (Theorem~\ref{thm:decomp}), an explicit non-degeneracy construction making the
identification conditions consistent with the generative model (Proposition~\ref{prop:law}),
a strict non-reducibility theorem (Theorem~\ref{thm:nonred}), and matched
quantitative/qualitative test batteries with their precise logical status
(Propositions~\ref{prop:quant}--\ref{prop:qual}). We are explicit throughout about which
claims are algebraic identities, which are causal claims requiring stronger assumptions, and
which are modeling hypotheses rather than theorems.
\end{abstract}

\section{Scope and what is and is not proved}

The substantive question --- how structural features and shared interpersonal beliefs jointly
produce team learning --- is a modeling task, not a theorem about the world. We therefore
\emph{separate} the modeling hypothesis (the directions of the arrows and the dual mediator/
moderator role of psychological safety) from the mathematical statements we can actually
prove about \emph{any} model with that form. Concretely:
\begin{itemize}
\item The choice of DAG and of which terms appear in each equation is a substantive
hypothesis. We do not (and cannot) prove it true of real teams; instead we
\emph{justify} it by showing it is the minimal linear recursive model consistent with a short
list of qualitative requirements extracted from the prompt (Theorem~\ref{thm:rep}), and we
make every implication of the model \emph{falsifiable} (\S\ref{sec:test}).
\item Statements labeled ``regression/projection'' are pure consequences of
orthogonality (Assumption~\ref{ass:exog}); they require no causal assumptions but
identify projections, \emph{not} causal effects.
\item Statements labeled ``causal'' (interventional means, natural/controlled effects)
require the additional Assumption~\ref{ass:ign} (sequential ignorability / back-door),
which we state explicitly and whose role we flag wherever it is used.
\end{itemize}

\section{Formalization}

\begin{definition}[Team variables]\label{def:vars}
Fix a population of work teams on a probability space $(\Omega,\mathcal F,\mathbb P)$. For a
team define team-level random quantities with finite second moments:
\begin{itemize}
\item $\bm X=(X_1,\dots,X_p)\in\mathbb R^p$: \emph{structural features} (team design, leader
coaching, context support, composition);
\item $S\in\mathbb R$: \emph{psychological safety}, the shared belief that the team is safe
for interpersonal risk-taking, operationalized as a team-level latent obtained by aggregating
member perceptions (Definition~\ref{def:shared});
\item $L\in\mathbb R$: \emph{learning behavior} (feedback seeking, information sharing, help
seeking, experimentation, error discussion), a team-level composite;
\item $P\in\mathbb R$: downstream \emph{team performance}.
\end{itemize}
\end{definition}

\begin{definition}[Sharedness / aggregation validity]\label{def:shared}
Let $S_{ij}$ be member $j$'s perception in team $i$, $j=1,\dots,n_i$. The construct $S$ is
\emph{shared} if within-team agreement is high and between-team variance dominates,
formalized by an intraclass correlation
$\mathrm{ICC}(1)=\sigma^2_{\text{between}}/(\sigma^2_{\text{between}}+\sigma^2_{\text{within}})$
bounded away from $0$; then the team mean $S_i=\frac1{n_i}\sum_j S_{ij}$ is a consistent
estimator of the team-level latent as $n_i\to\infty$, licensing treatment of $S$ as one
team-level variable. This is a measurement-model precondition, assumed throughout.
\end{definition}

\section{The integrated structural model and its justification}

\begin{model}[Integrated structural--interpersonal SCM]\label{mod:main}
The framework asserts the recursive (acyclic) generative process
\begin{align}
S &= \alpha_0 + \bm\alpha^\top \bm X + \varepsilon_S, \label{eq:S}\\
L &= \beta_0 + \bm\beta^\top \bm X + \gamma\, S
      + (\bm\delta^\top \bm X)\, S + \varepsilon_L, \label{eq:L}\\
P &= \eta_0 + \theta\, L + \bm\phi^\top \bm X + \varepsilon_P, \label{eq:P}
\end{align}
with topological order $\bm X \prec S \prec L \prec P$ on the DAG with edges
$\bm X\to S$, $\bm X\to L$, $S\to L$, $L\to P$, $\bm X\to P$, and the moderating mechanism
encoded by the product regressor $(\bm\delta^\top\bm X)S$. Here
$\bm\alpha,\bm\beta,\bm\delta,\bm\phi\in\mathbb R^p$,
$\alpha_0,\beta_0,\eta_0,\gamma,\theta\in\mathbb R$, and $(\varepsilon_S,\varepsilon_L,
\varepsilon_P)$ are mean-zero disturbances. The slope of $L$ on $S$ is
$\gamma+\bm\delta^\top\bm X$, i.e.\ structure conditions the safety--learning link.
\end{model}

The two substantive roles are read off the equations: in \eqref{eq:S} structure is an
\emph{antecedent} of safety; in \eqref{eq:L} safety \emph{mediates} part of the structure
effect and, through $(\bm\delta^\top\bm X)S$, also \emph{moderates} the safety--learning link
($\partial^2 \mathbb E[L\mid\bm X,S]/\partial\bm X\,\partial S=\bm\delta$). The following
theorem upgrades ``read off'' to a justification: the model is the canonical minimal one
meeting explicit qualitative requirements.

\begin{theorem}[Representation / minimality]\label{thm:rep}
Consider the class $\mathcal M$ of recursive structural models for $(\bm X,S,L,P)$ that
satisfy all of the following qualitative desiderata, taken directly from the prompt:
\begin{enumerate}
\item[(D1)] \emph{Topological order} $\bm X\prec S\prec L\prec P$ (structure is exogenous;
safety is an interpersonal belief formed given structure; learning is behavior given the
belief and structure; performance is downstream of learning);
\item[(D2)] each structural equation's systematic part is an \emph{affine function of its
parents plus at most pairwise products of parents} (linearity with interactions; the
smallest polynomial class admitting moderation);
\item[(D3)] \emph{antecedence}: $S$ may depend on $\bm X$ (an $S$ equation with an
$\bm X$ term is permitted);
\item[(D4)] \emph{mediation}: $\bm X$ may affect $L$ both directly and through $S$ (both an
$\bm X$ term and an $S$ term in the $L$ equation);
\item[(D5)] \emph{moderation of the safety--learning link by structure}: the partial
derivative $\partial \mathbb E[L\mid\bm X,S]/\partial S$ is permitted to depend on $\bm X$;
\item[(D6)] \emph{no moderation that the prompt does not name}: the only product term
retained is the one required by (D5); in particular the $S$ equation has no $\bm X$--$\bm X$
products beyond linearity, and $P$ depends on $\{\bm X,L\}$ affinely.
\end{enumerate}
Then the systematic (mean) part of every model in $\mathcal M$ has exactly the functional
form of \eqref{eq:S}--\eqref{eq:P}, and \eqref{eq:L} is the unique form in $\mathcal M$ whose
$S$-slope depends on $\bm X$; dropping the term $(\bm\delta^\top\bm X)S$ violates (D5),
while adding any further pairwise product permitted by (D2) violates (D6). Hence
Model~\ref{mod:main} is the minimal element of $\mathcal M$ realizing all of (D1)--(D5).
\end{theorem}

\begin{proof}
By (D1) the equations are recursive in the order $\bm X\prec S\prec L\prec P$, so the
systematic part of $S$ is a function of $\bm X$ only, that of $L$ of $(\bm X,S)$ only, and
that of $P$ of $(\bm X,L)$ only. By (D2) each is affine in its parents plus pairwise
products of parents.

\emph{$S$ equation.} Parents are $\bm X$; the affine-plus-pairwise class is
$\alpha_0+\bm\alpha^\top\bm X+\bm X^\top A\bm X$ for some symmetric $A$. (D6) forbids
$\bm X$--$\bm X$ products in the $S$ equation (the prompt names no such interaction), so
$A=0$ and we recover \eqref{eq:S}.

\emph{$L$ equation.} Parents are $(\bm X,S)$. The affine-plus-pairwise class is
$\beta_0+\bm\beta^\top\bm X+\gamma S+\bm X^\top B\bm X+(\bm\delta^\top\bm X)S+ cS^2$. The
$S$-slope is $\gamma+\bm\delta^\top\bm X+2cS+\dots$; (D5) requires this slope to depend on
$\bm X$, which forces $\bm\delta\neq0$ to be \emph{permitted}, i.e.\ the term
$(\bm\delta^\top\bm X)S$ must be retained. (D6) retains \emph{only} the product needed for
(D5): the $\bm X$--$S$ product. Hence $B=0$, $c=0$, and we recover \eqref{eq:L}. If
$(\bm\delta^\top\bm X)S$ is dropped ($\bm\delta\equiv0$ identically as a constraint) the
$S$-slope is the constant $\gamma$ and (D5) fails; if $\bm X^\top B\bm X$ or $cS^2$ is added,
(D6) fails. Thus \eqref{eq:L} is the unique minimal form.

\emph{$P$ equation.} Parents are $(\bm X,L)$; (D6) restricts $P$ to be affine in them,
giving \eqref{eq:P}.

Existence and acyclicity of the disturbances follow by appending mean-zero residuals
$\varepsilon_S,\varepsilon_L,\varepsilon_P$, which is always possible. Minimality is now the
conjunction of the three paragraphs.
\end{proof}

\begin{remark}
Theorem~\ref{thm:rep} is a justification \emph{relative to the desiderata} (D1)--(D6); those
desiderata are the substantive content imported from the team-learning prompt and are not
themselves theorems. The contribution is to make the modeling choice canonical and explicit
rather than ad hoc: any other functional form either drops a named mechanism or adds an
unnamed one.
\end{remark}

\section{Assumptions: regression vs.\ causal}

\begin{assumption}[Recursive exogeneity / disturbance orthogonality]\label{ass:exog}
Conditional on the regressors appearing in each equation, the disturbances are mean
independent and mutually orthogonal across equations:
\[
\mathbb E[\varepsilon_S\mid \bm X]=0,\quad
\mathbb E[\varepsilon_L\mid \bm X,S]=0,\quad
\mathbb E[\varepsilon_P\mid \bm X,L]=0,
\]
and $\mathbb E[\varepsilon_S\varepsilon_L]=\mathbb E[\varepsilon_S\varepsilon_P]
=\mathbb E[\varepsilon_L\varepsilon_P]=0$. Since $(\bm\delta^\top\bm X)S$ is a measurable
function of $(\bm X,S)$, the middle display also gives
$\mathbb E[\varepsilon_L\mid \bm X,S,(\bm\delta^\top\bm X)S]=0$.
\end{assumption}

Assumption~\ref{ass:exog} is purely about the joint law of the disturbances and regressors;
it suffices for projection identification (Theorem~\ref{thm:ident}) but \emph{not} for any
causal reading. For causal claims we add:

\begin{assumption}[Sequential ignorability / back-door]\label{ass:ign}
Using potential-outcome notation, let $S(\bm x)$ be the potential value of safety under an
intervention setting $\bm X=\bm x$, and $L(\bm x,s)$ the potential learning behavior under
$\bm X=\bm x$, $S=s$. We assume:
\begin{enumerate}
\item[(I1)] \emph{Exposure ignorability}: $\bm X$ is (conditionally) randomized, e.g.\ by
design (randomized coaching/support) or by an observed sufficient adjustment set; here we take
$\bm X$ exogenous, $\{S(\bm x),L(\bm x,s)\}\perp \bm X$.
\item[(I2)] \emph{Mediator ignorability (no mediator--outcome confounding)}:
$L(\bm x,s)\perp S\mid \bm X$, i.e.\ given structure there is no unobserved cause of both
safety and learning. This is the Imai--Keele--Yamamoto sequential-ignorability condition for
a single mediator.
\item[(I3)] \emph{Consistency/SUTVA} and positivity (overlap) of $\bm X,S$.
\end{enumerate}
\end{assumption}

We flag explicitly: (I2) is a strong, non-testable identifying assumption; it is exactly
what licenses interpreting the mediation/moderation algebra causally, and we never use it
where it is not needed.

\section{Causal layer: interventional means equal conditional means}

\begin{theorem}[do-means$=$conditional means]\label{thm:causal}
Under Model~\ref{mod:main}, Assumption~\ref{ass:exog} and Assumption~\ref{ass:ign},
\begin{align}
\mathbb E\big[L(\bm x,s)\big]&=\mathbb E[L\mid \bm X=\bm x,S=s]
   =\beta_0+\bm\beta^\top\bm x+(\gamma+\bm\delta^\top\bm x)s,\label{eq:doLxs}\\
\mathbb E\big[S(\bm x)\big]&=\mathbb E[S\mid \bm X=\bm x]=\alpha_0+\bm\alpha^\top\bm x,
\label{eq:doSx}\\
\mathbb E\big[L(\bm x,S(\bm x'))\big]
  &=\beta_0+\bm\beta^\top\bm x+(\gamma+\bm\delta^\top\bm x)(\alpha_0+\bm\alpha^\top\bm x').
\label{eq:nested}
\end{align}
\end{theorem}

\begin{proof}
By (I1) and consistency, $\mathbb E[S(\bm x)]=\mathbb E[S(\bm x)\mid\bm X=\bm x]
=\mathbb E[S\mid\bm X=\bm x]$, and \eqref{eq:S} with Assumption~\ref{ass:exog} gives
$\mathbb E[S\mid\bm X=\bm x]=\alpha_0+\bm\alpha^\top\bm x$, proving \eqref{eq:doSx}.
By (I1)--(I2) and consistency, $\mathbb E[L(\bm x,s)]=\mathbb E[L(\bm x,s)\mid\bm X=\bm x,
S=s]=\mathbb E[L\mid\bm X=\bm x,S=s]$; the last equals
$\beta_0+\bm\beta^\top\bm x+(\gamma+\bm\delta^\top\bm x)s$ by \eqref{eq:L} and
$\mathbb E[\varepsilon_L\mid\bm X,S]=0$, proving \eqref{eq:doLxs}. For the nested
counterfactual, by (I2) ($L(\bm x,\cdot)\perp S\mid\bm X$) and (I1),
\[
\mathbb E[L(\bm x,S(\bm x'))]
=\mathbb E_{s\sim S(\bm x')}\mathbb E[L(\bm x,s)]
=\mathbb E[L(\bm x,s)]\big|_{s=\mathbb E[S(\bm x')]},
\]
the inner mean being affine in $s$; substituting \eqref{eq:doLxs} and \eqref{eq:doSx}
yields \eqref{eq:nested}.
\end{proof}

\section{Exact four-way effect decomposition}

We now give the \emph{canonical} causal decomposition (VanderWeele's four-way
decomposition) of the total effect of moving structure from $\bm x'$ to $\bm x$ on learning,
with each term carrying its standard meaning. Fix a reference safety level $s^\ast$ (commonly
$s^\ast=0$ after centering). Define, under Assumption~\ref{ass:ign}, the total effect
$\mathrm{TE}=\mathbb E[L(\bm x,S(\bm x))]-\mathbb E[L(\bm x',S(\bm x'))]$, and write
$m(\bm x):=\alpha_0+\bm\alpha^\top\bm x=\mathbb E[S(\bm x)]$.

\begin{theorem}[Four-way decomposition]\label{thm:decomp}
Under Model~\ref{mod:main}, Assumption~\ref{ass:exog} and Assumption~\ref{ass:ign}, for any
$\bm x,\bm x'\in\mathbb R^p$ and reference $s^\ast\in\mathbb R$,
\begin{equation}\label{eq:4way}
\mathrm{TE}
=\underbrace{(\bm\beta+\bm\delta\,s^\ast)^\top(\bm x-\bm x')}_{\mathrm{CDE}(s^\ast)}
+\underbrace{\bm\delta^\top(\bm x-\bm x')\,(m(\bm x')-s^\ast)}_{\mathrm{INT_{ref}}}
+\underbrace{\bm\delta^\top(\bm x-\bm x')\,(m(\bm x)-m(\bm x'))}_{\mathrm{INT_{med}}}
+\underbrace{(\gamma+\bm\delta^\top\bm x')\,(m(\bm x)-m(\bm x'))}_{\mathrm{PIE}},
\end{equation}
where $\mathrm{CDE}(s^\ast)$ is the controlled direct effect at $S=s^\ast$,
$\mathrm{INT_{ref}}$ the reference (pure) interaction, $\mathrm{INT_{med}}$ the mediated
interaction, and $\mathrm{PIE}$ the pure indirect effect; these are exactly the four terms of
VanderWeele's decomposition for an exposure with a single mediator and exposure--mediator
interaction. Consequently the \emph{natural direct} and \emph{natural indirect} effects are
\begin{align}
\mathrm{NDE}&=\mathbb E[L(\bm x,S(\bm x'))]-\mathbb E[L(\bm x',S(\bm x'))]
   =\mathrm{CDE}(s^\ast)+\mathrm{INT_{ref}}
   =(\bm\beta+\bm\delta\, m(\bm x'))^\top(\bm x-\bm x'),\\
\mathrm{NIE}&=\mathbb E[L(\bm x,S(\bm x))]-\mathbb E[L(\bm x,S(\bm x'))]
   =\mathrm{INT_{med}}+\mathrm{PIE}
   =(\gamma+\bm\delta^\top\bm x)\,(m(\bm x)-m(\bm x')),
\end{align}
and $\mathrm{TE}=\mathrm{NDE}+\mathrm{NIE}$. When $\bm\delta=\bm0$ (no moderation), all
interaction terms vanish and $\mathrm{TE}=\bm\beta^\top(\bm x-\bm x')+\gamma\,
\bm\alpha^\top(\bm x-\bm x')=(\bm\beta+\gamma\bm\alpha)^\top(\bm x-\bm x')$, the classical
linear mediation identity \emph{total $=$ direct $+$ indirect}.
\end{theorem}

\begin{proof}
By \eqref{eq:nested}, $\mathbb E[L(\bm x,S(\bm x))]
=\beta_0+\bm\beta^\top\bm x+(\gamma+\bm\delta^\top\bm x)m(\bm x)$ and likewise at
$\bm x'$, so
\[
\mathrm{TE}=\bm\beta^\top(\bm x-\bm x')
+(\gamma+\bm\delta^\top\bm x)m(\bm x)-(\gamma+\bm\delta^\top\bm x')m(\bm x').
\]
We verify \eqref{eq:4way} by expanding its four terms and summing:
\begin{align*}
\mathrm{CDE}+\mathrm{INT_{ref}}
&=(\bm\beta+\bm\delta s^\ast)^\top(\bm x-\bm x')
  +\bm\delta^\top(\bm x-\bm x')(m(\bm x')-s^\ast)\\
&=\bm\beta^\top(\bm x-\bm x')+\bm\delta^\top(\bm x-\bm x')\,m(\bm x')
  =\mathrm{NDE},
\end{align*}
(the $s^\ast$ terms cancel, so the decomposition is independent of $s^\ast$ up to the split
between CDE and INT$_{\mathrm{ref}}$). Next,
\begin{align*}
\mathrm{INT_{med}}+\mathrm{PIE}
&=\bm\delta^\top(\bm x-\bm x')(m(\bm x)-m(\bm x'))
  +(\gamma+\bm\delta^\top\bm x')(m(\bm x)-m(\bm x'))\\
&=(m(\bm x)-m(\bm x'))\big[\bm\delta^\top\bm x-\bm\delta^\top\bm x'+\gamma+\bm\delta^\top\bm
x'\big]
=(\gamma+\bm\delta^\top\bm x)(m(\bm x)-m(\bm x'))=\mathrm{NIE}.
\end{align*}
Adding,
\begin{align*}
\mathrm{NDE}+\mathrm{NIE}
&=\bm\beta^\top(\bm x-\bm x')+\bm\delta^\top(\bm x-\bm x')m(\bm x')
  +(\gamma+\bm\delta^\top\bm x)(m(\bm x)-m(\bm x'))\\
&=\bm\beta^\top(\bm x-\bm x')+(\gamma+\bm\delta^\top\bm x)m(\bm x)
  -(\gamma+\bm\delta^\top\bm x')m(\bm x')\\
&\qquad+\big[\bm\delta^\top(\bm x-\bm x')m(\bm x')-(\gamma+\bm\delta^\top\bm x)m(\bm x')
  +(\gamma+\bm\delta^\top\bm x')m(\bm x')\big].
\end{align*}
The bracket equals $m(\bm x')[\bm\delta^\top\bm x-\bm\delta^\top\bm x'-\gamma-\bm\delta^\top
\bm x+\gamma+\bm\delta^\top\bm x']=0$, leaving exactly $\mathrm{TE}$. The $\bm\delta=\bm0$
specialization is immediate since then $m(\bm x)-m(\bm x')=\bm\alpha^\top(\bm x-\bm x')$.
\end{proof}

\begin{remark}[On the previous, non-canonical labels]
A naive regrouping of the bilinear terms of $(\gamma+\bm\delta^\top\bm x)m(\bm x)$ into
``direct $=\bm\beta^\top(\bm x-\bm x')$,'' ``mediated,'' and ``moderated baseline''
\emph{also} sums to $\mathrm{TE}$, but those labels do \emph{not} coincide with the
canonical natural/controlled effects except when $\bm x'=\bm0$: e.g.\ at $\bm
x'=(-0.3,0.2),\ \bm x=(1,0.5)$ with $\alpha_0=.5,\bm\alpha=(.3,-.2),\bm\beta=(.4,.1),
\gamma=.7,\bm\delta=(.5,-.3)$ one has $\mathrm{NDE}=0.757\neq\bm\beta^\top(\bm x-\bm x')=
0.55$. We therefore use only the VanderWeele terms above, whose causal meaning is standard
and reference-consistent.
\end{remark}

\begin{corollary}[Performance pathway]\label{cor:perf}
Under Model~\ref{mod:main} and Assumptions~\ref{ass:exog},~\ref{ass:ign} (with the analogous
ignorability for the $L\to P$ link, $P(\ell)\perp L\mid\bm X$),
\[
\mathbb E[P(\bm x)]=\eta_0+\bm\phi^\top\bm x+\theta\,\mathbb E[L(\bm x,S(\bm x))],
\]
so the total causal effect of structure on performance equals the direct structural part
$\bm\phi^\top(\bm x-\bm x')$ plus $\theta\cdot\mathrm{TE}$ with $\mathrm{TE}$ from
Theorem~\ref{thm:decomp}.
\end{corollary}
\begin{proof}
By the $L\to P$ ignorability and consistency, $\mathbb E[P\mid\bm X=\bm x,L=\ell]
=\eta_0+\bm\phi^\top\bm x+\theta\ell$ (using \eqref{eq:P} and
$\mathbb E[\varepsilon_P\mid\bm X,L]=0$) equals $\mathbb E[P(\ell)\mid\bm X=\bm x]$; averaging
over $L(\bm x,S(\bm x))$ and using affineness in $\ell$ gives the display. Subtracting at
$\bm x'$ and applying Theorem~\ref{thm:decomp} gives the effect statement.
\end{proof}

\section{Statistical identification (regression layer)}

The decomposition's ingredients $(\bm\alpha,\bm\beta,\gamma,\bm\delta,\theta,\bm\phi)$ are
recoverable from the joint law under Assumption~\ref{ass:exog} \emph{alone} (no causal
assumption needed to identify the coefficients; causal meaning is then attached via
Theorem~\ref{thm:causal}).

\begin{theorem}[Parameter identification]\label{thm:ident}
Suppose, in addition to Assumption~\ref{ass:exog}:
\begin{enumerate}
\item[(R1)] $\Sigma_Z:=\mathbb E[\bm Z\bm Z^\top]\succ0$ for
$\bm Z:=(1,\bm X^\top,S,(\bm X S)^\top)^\top\in\mathbb R^{2p+2}$ (the regressors of
\eqref{eq:L});
\item[(R2)] the second-moment matrix of $(1,\bm X^\top)$ is nonsingular;
\item[(R3)] the second-moment matrix of $(1,\bm X^\top,L)$ is nonsingular.
\end{enumerate}
Then $(\alpha_0,\bm\alpha)$, $(\beta_0,\bm\beta,\gamma,\bm\delta)$, and
$(\eta_0,\bm\phi,\theta)$ are uniquely determined by the joint distribution of
$(\bm X,S,L,P)$, each as a population least-squares projection; consequently every term in
\eqref{eq:4way} and Corollary~\ref{cor:perf} is identified.
\end{theorem}

\begin{proof}
\emph{Eq.\ \eqref{eq:S}.} $\mathbb E[\varepsilon_S\mid\bm X]=0$ gives
$\mathbb E[(1,\bm X^\top)^\top\varepsilon_S]=\bm0$, so $(\alpha_0,\bm\alpha)$ solve the normal
equations $\mathbb E[(1,\bm X^\top)^\top(1,\bm X^\top)](\alpha_0,\bm\alpha^\top)^\top
=\mathbb E[(1,\bm X^\top)^\top S]$; by (R2) the matrix is invertible, so the solution is the
unique projection of $S$ on $(1,\bm X)$.

\emph{Eq.\ \eqref{eq:L}.} Each component of $\bm Z$ is $(\bm X,S)$-measurable and
$\mathbb E[\varepsilon_L\mid\bm X,S]=0$, so $\mathbb E[\bm Z\,\varepsilon_L]=\bm0$ and the
coefficient vector $\bm b:=(\beta_0,\bm\beta^\top,\gamma,\bm\delta^\top)^\top$ solves
$\Sigma_Z\bm b=\mathbb E[\bm Z\,L]$; by (R1), $\bm b=\Sigma_Z^{-1}\mathbb E[\bm Z\,L]$ is
unique. In particular $\bm\delta$ is identified \emph{whenever (R1) holds}.

\emph{Eq.\ \eqref{eq:P}.} $\mathbb E[\varepsilon_P\mid\bm X,L]=0$ gives
$\mathbb E[(1,\bm X^\top,L)^\top\varepsilon_P]=\bm0$; by (R3) the second-moment matrix is
invertible and $(\eta_0,\bm\phi,\theta)$ is the unique projection.

Every parameter is a fixed matrix functional of population moments; \eqref{eq:4way} and
Corollary~\ref{cor:perf} are continuous in those parameters and in $\bm x,\bm x'$, hence
identified.
\end{proof}

The next proposition resolves a subtlety the critic raised: (R1) must be \emph{compatible}
with the generative model, since under \eqref{eq:S} the regressors $S$ and $\bm X S$ are
functions of $\bm X$. We exhibit an explicit admissible law making $\Sigma_Z\succ0$, and we
flag the degenerate case where identification of $\bm\delta$ fails.

\begin{proposition}[Admissible law; degenerate exception]\label{prop:law}
Let $\bm X\sim N(\bm0,\Sigma_X)$ with $\Sigma_X\succ0$ and let
$\varepsilon_S\sim N(0,\sigma_S^2)$ be independent of $\bm X$ with $\sigma_S^2>0$, and set
$S=\alpha_0+\bm\alpha^\top\bm X+\varepsilon_S$. Then $\Sigma_Z\succ0$, so (R1) holds and all
coefficients (including $\bm\delta$) are identified. Conversely, if $\sigma_S^2=0$ (perfect
mediation, $\varepsilon_S\equiv0$), then $S$ and $\bm X S$ are polynomials in $\bm X$, the
vector $\bm Z$ lies in a proper subspace of $\mathbb R^{2p+2}$, $\Sigma_Z$ is singular, and
$\bm\delta$ is \emph{not} identified by \eqref{eq:L} alone; thus the claim that moderation is
``empirically decidable'' holds precisely when $\sigma_S^2>0$ (and more generally when (R1)
holds).
\end{proposition}

\begin{proof}
\emph{Nonsingularity.} Suppose $\bm c^\top\bm Z=0$ a.s.\ for some
$\bm c=(c_0,\bm c_1^\top,c_2,\bm c_3^\top)^\top$; we show $\bm c=\bm0$. Write
$\bm c^\top\bm Z=c_0+\bm c_1^\top\bm X+c_2 S+(\bm c_3^\top\bm X)S=0$ a.s. Substituting
$S=\alpha_0+\bm\alpha^\top\bm X+\varepsilon_S$ and grouping by powers of $\varepsilon_S$:
the coefficient of $\varepsilon_S$ is $c_2+\bm c_3^\top\bm X$, which must vanish a.s.\ as a
function of $\bm X$ (since, conditional on $\bm X$, $\varepsilon_S$ has positive variance);
as $\Sigma_X\succ0$, the affine function $c_2+\bm c_3^\top\bm X\equiv0$ forces $c_2=0$ and
$\bm c_3=\bm0$. The remaining $\varepsilon_S^0$ part is
$c_0+\alpha_0 c_2+(\bm c_1+c_2\bm\alpha)^\top\bm X=c_0+\bm c_1^\top\bm X\equiv0$, which by
$\Sigma_X\succ0$ forces $c_0=0,\bm c_1=\bm0$. Hence $\bm c=\bm0$ and $\Sigma_Z\succ0$.

\emph{Degeneracy.} If $\sigma_S^2=0$ then $S=\alpha_0+\bm\alpha^\top\bm X$ and
$\bm X S=\alpha_0\bm X+\bm X(\bm\alpha^\top\bm X)$, each a (degree $\le2$) polynomial in
$\bm X$. Take $\bm c$ with $c_2=1,\ c_0=-\alpha_0,\ \bm c_1=-\bm\alpha,\ \bm c_3=\bm0$: then
$\bm c^\top\bm Z=-\alpha_0-\bm\alpha^\top\bm X+S=\varepsilon_S=0$ a.s., a nontrivial linear
dependence, so $\Sigma_Z$ is singular. Then the normal equations $\Sigma_Z\bm b
=\mathbb E[\bm Z L]$ have a nontrivial null space, so $\bm b$ (in particular $\bm\delta$) is
not uniquely determined by \eqref{eq:L} alone.
\end{proof}

\begin{remark}[Estimation and inference]
Theorem~\ref{thm:ident} is constructive: each parameter is a population least-squares
projection. Replacing population by sample moments yields the OLS/SEM estimators, consistent
and asymptotically normal under (R1)--(R3) and finite fourth moments; all effects in
\eqref{eq:4way} then admit delta-method or bootstrap confidence intervals. In the
no-moderation case $\gamma\bm\alpha$ is the classical product-of-coefficients mediation
estimand.
\end{remark}

\section{Non-reducibility to independent predictors}

\begin{theorem}[Strict superiority of the joint model]\label{thm:nonred}
Let the data be generated by Model~\ref{mod:main} with $\bm\delta\neq\bm0$ and with the
interaction regressor not collinear with the additive ones (the $L^2$-residual of
$(\bm\delta^\top\bm X)S$ after projecting on $\overline{\mathrm{span}}\{1,\bm X,S\}$ is
nonzero; this holds under (R1)). Let
$g_{\mathrm{add}}(\bm X,S)=\mu_0+\bm\mu^\top\bm X+\nu S$ range over all affine
(no-interaction) predictors and
$g_{\mathrm{int}}(\bm X,S)=\beta_0+\bm\beta^\top\bm X+\gamma S+(\bm\delta^\top\bm X)S$ be the
conditional mean of $L$ under \eqref{eq:L}. Then
\[
\min_{\mu_0,\bm\mu,\nu}\mathbb E[(L-g_{\mathrm{add}})^2]
>\mathbb E[(L-g_{\mathrm{int}})^2],
\]
the gap equal to $\sigma_R^2:=\mathbb E[R^2]>0$, the variance of the $L^2$-projection
residual $R$ of $(\bm\delta^\top\bm X)S$.
\end{theorem}

\begin{proof}
Let $W:=(\bm\delta^\top\bm X)S$, $\mathcal A:=\overline{\mathrm{span}}\{1,X_1,\dots,X_p,S\}$,
$\hat W:=\mathrm{proj}_{\mathcal A}W$, $R:=W-\hat W\perp\mathcal A$ with
$\mathbb E[R^2]=\sigma_R^2>0$. By \eqref{eq:L},
$L=g_{\mathrm{int}}(\bm X,S)+\varepsilon_L$ with $\varepsilon_L\perp\mathcal A$ and
$\varepsilon_L\perp W$ (as $W$ is $(\bm X,S)$-measurable), hence $\varepsilon_L\perp R$.
Write $L=(\beta_0+\bm\beta^\top\bm X+\gamma S+\hat W)+R+\varepsilon_L$ with the first bracket
in $\mathcal A$ and $R,\varepsilon_L\perp\mathcal A$. So the projection of $L$ onto
$\mathcal A$ is the bracket, and
$\min_{\mathcal A}\mathbb E[(L-g_{\mathrm{add}})^2]=\mathbb E[(R+\varepsilon_L)^2]
=\sigma_R^2+\mathbb E[\varepsilon_L^2]$ (by $R\perp\varepsilon_L$). Since
$\mathbb E[(L-g_{\mathrm{int}})^2]=\mathbb E[\varepsilon_L^2]$, the difference is
$\sigma_R^2>0$.
\end{proof}

\begin{corollary}[Bias of the ``independent predictors'' slope]\label{cor:bias}
Under the hypotheses of Theorem~\ref{thm:nonred}, the additive least-squares slope on $S$,
$\nu^\ast$, generally differs from $\gamma$ and from $\gamma+\bm\delta^\top\bm x$ at any
$\bm x$. Hence treating safety and structure as independent additive predictors both loses
fit (Theorem~\ref{thm:nonred}) and yields a configuration-blind, biased estimate of the
safety--learning relationship.
\end{corollary}
\begin{proof}
$\nu^\ast$ is the $S$-coefficient of $\mathrm{proj}_{\mathcal A}L$, which loads on $S$
through $\hat W$ whenever $(\bm\delta^\top\bm X)S$ has nonzero partial covariance with $S$
given $(1,\bm X)$, so $\nu^\ast\neq\gamma$ in general. The true marginal effect of $S$ at
$\bm X=\bm x$ is $\partial\mathbb E[L\mid\bm X=\bm x,S=s]/\partial s=\gamma+\bm\delta^\top
\bm x$, varying with $\bm x$, hence not equal to the single constant $\nu^\ast$ unless
$\bm\delta=\bm0$.
\end{proof}

\section{Dual quantitative--qualitative testability}\label{sec:test}

\begin{proposition}[Quantitative test]\label{prop:quant}
Under Assumption~\ref{ass:exog} and (R1)--(R3), each of the following is point-identified and
estimable by least squares/SEM, with a falsifiable null:
\begin{enumerate}
\item \emph{Antecedent} $\bm\alpha\neq\bm0$: test $H_0:\bm\alpha=\bm0$ in \eqref{eq:S}.
\item \emph{Mediation} $\gamma\bm\alpha\neq\bm0$: product-of-coefficients/bootstrap test of
$H_0:\gamma\bm\alpha=\bm0$ (causal meaning under Assumption~\ref{ass:ign}).
\item \emph{Moderation} $\bm\delta\neq\bm0$: test the interaction block in \eqref{eq:L}; by
Theorem~\ref{thm:nonred} (and Proposition~\ref{prop:law} for identifiability) $\bm\delta\neq
\bm0$ iff the integrated model strictly improves population fit over the additive one.
\item \emph{Performance channel} $\theta\neq0$ and the mediated $L$-effect on $P$
(Corollary~\ref{cor:perf}).
\end{enumerate}
\end{proposition}
\begin{proof}
Each parameter is identified (Theorem~\ref{thm:ident}, Proposition~\ref{prop:law}) and
consistently estimable as a least-squares projection (Remark after Theorem~\ref{thm:ident});
a parameter restriction is a smooth hypothesis, falsifiable by Wald/LR/bootstrap. Item~3's
equivalence is Theorem~\ref{thm:nonred}.
\end{proof}

\begin{proposition}[Qualitative / mechanism test]\label{prop:qual}
The model makes the following \emph{exactly} stated, case-checkable predictions.
\begin{enumerate}
\item \emph{Interventional mediation (not a temporal claim about a static model).} Under
Assumption~\ref{ass:ign}, an exogenous structural change $\bm x'\!\to\!\bm x$ produces an
expected safety change $m(\bm x)-m(\bm x')=\bm\alpha^\top(\bm x-\bm x')$ and an expected
indirect learning change $\mathrm{NIE}=(\gamma+\bm\delta^\top\bm x)(m(\bm x)-m(\bm x'))$
(Theorem~\ref{thm:decomp}); a case study that intervenes on structure should see safety move
by the predicted amount and learning move accordingly through that channel. (If a
\emph{temporal} ordering ``$\bm X$-change precedes $S$-change precedes $L$-change'' is to be
asserted, the static SCM must be replaced by a longitudinal SCM with time-indexed variables
$S_t,L_t$; the cross-sectional model entails the interventional, not the temporal, claim.)
\item \emph{Moderation contrast (exact).} For two teams at structures $\bm x_1,\bm x_2$ with
equal $S=s$, the model gives $\mathbb E[L\mid\bm X=\bm x_1,S=s]-\mathbb E[L\mid\bm X=\bm
x_2,S=s]=\bm\beta^\top(\bm x_1-\bm x_2)+s\,\bm\delta^\top(\bm x_1-\bm x_2)$, so the
\emph{same} safety yields a structure-dependent amount of learning, with slope difference
$\bm\delta^\top(\bm x_1-\bm x_2)$ --- a contrast observable in comparative case study.
\item \emph{Exact conditional independence (iff).} In a recursive linear SCM the global
Markov property holds, so $d$-separation gives: $P\perp\bm X\mid L$ \emph{if and only if}
$\bm\phi=\bm0$ (and the disturbances satisfy Assumption~\ref{ass:exog}); and conditional on
its parents and their product $(\bm X,S,\bm X S)$, $L$ depends on $\bm X$ only through those
regressors, i.e.\ $\mathbb E[L\mid\bm X,S]$ is the stated function with no residual
$\bm X$-dependence. Observed violations falsify the corresponding edge restriction.
\end{enumerate}
\end{proposition}
\begin{proof}
Item~1: the displayed quantities are Theorem~\ref{thm:decomp}/\eqref{eq:doSx}; the
parenthetical is the standard fact that a static SCM fixes only a logical ordering, so the
temporal claim requires time-indexed variables. Item~2: subtract \eqref{eq:doLxs} at
$(\bm x_1,s)$ and $(\bm x_2,s)$. Item~3: by the global Markov property of the recursive
linear SCM, $P$ and $\bm X$ are $d$-separated by $L$ exactly when the direct edge
$\bm X\to P$ is absent, i.e.\ $\bm\phi=\bm0$; the ``only if'' direction is because a nonzero
$\bm\phi$ creates an active path $\bm X\to P$ not blocked by conditioning on $L$, yielding
$\mathrm{Cov}(P,\bm X\mid L)\neq0$ generically. The residual-independence statement for $L$
is immediate from \eqref{eq:L} since, given $(\bm X,S,\bm X S)$, $L=\,$(stated mean)$+
\varepsilon_L$ with $\mathbb E[\varepsilon_L\mid\bm X,S]=0$.
\end{proof}

\section{Conclusion}

Model~\ref{mod:main} is the requested integrated framework, and we have separated and proved
its three layers. Theorem~\ref{thm:rep} \emph{justifies} the functional form as the minimal
linear recursive model meeting explicit qualitative desiderata extracted from the prompt.
Under disturbance orthogonality (Assumption~\ref{ass:exog}) all coefficients are identified
as projections (Theorem~\ref{thm:ident}), with an explicit admissible law and a flagged
degenerate exception (Proposition~\ref{prop:law}). Under the additional, explicitly stated
sequential-ignorability assumption (Assumption~\ref{ass:ign}) the conditional means are
interventional means (Theorem~\ref{thm:causal}), so psychological safety enters \emph{both}
as a mediator and as a moderator, and the total effect decomposes \emph{exactly} into the
canonical four VanderWeele terms (Theorem~\ref{thm:decomp}), with performance driven through
learning (Corollary~\ref{cor:perf}). Theorem~\ref{thm:nonred} and Corollary~\ref{cor:bias}
prove the joint model strictly dominates any additive ``independent predictors'' treatment
whenever moderation is present, and Propositions~\ref{prop:quant}--\ref{prop:qual} provide
matched, precisely-scoped quantitative and qualitative test batteries. We have been explicit
about which statements are algebraic identities, which are causal claims contingent on
ignorability, and which are modeling hypotheses rather than theorems.
\end{document}
